\chapter{System Design and Implementation}
This chapter bridges the theoretical algorithms established previously with their practical application, providing a comprehensive technical overview of the software architecture engineered to natively support multi-fiber parallel links. It systematically details the structural enhancements integrated into the TeraFlowSDN ecosystem, the development of the underlying Routing and Spectrum Allocation (RSA) computational modules, and the operational deployment of the novel \texttt{ParallelOpticalController} microservice.
\section{Modifications to TeraFlowSDN}
To accommodate the complex spatial and spectral constraints of multi-fiber topologies, several core components of the TeraFlowSDN control plane required structural modifications. These enhancements ensure that the central database maintains a highly accurate, normalized representation of the physical optical layer.
\subsection{Enhancement of the DeviceService}
To accommodate the complex spatial and spectral constraints of multi-fiber topologies, several core components of the TeraFlowSDN control plane required structural modifications. These enhancements ensure that the central database maintains a highly accurate, normalized representation of the physical optical layer.
Modern elastic optical networks are inherently heterogeneous, relying on equipment from various manufacturers. While the industry utilizes standard YANG data models for device configuration and telemetry, the expansive scope of YANG allows different vendors to adopt significantly different internal data structures and naming conventions. Consequently, relying on a monolithic data extraction technique to parse equipment capabilities is unreliable and often fails to capture the granular optical parameters required for advanced RSA computation. To overcome this interoperability challenge, this thesis introduces a novel, dynamically adaptable data extraction technique integrated directly into the \texttt{DeviceService} component of TeraFlowSDN, which is responsible for orchestrating device drivers. Rather than applying a uniform parsing rule, the augmented \texttt{DeviceService} identifies the specific vendor and hardware model of the connected equipment and applies a tailored extraction profile. To maintain system robustness and ensure backward compatibility with uncharacterized or generic optical equipment, a default extraction methodology is enforced whenever a specific vendor signature is not recognized.

Following the initial data extraction, the driver logic applies a rigorous normalization phase. A standardized naming convention is forcefully applied to all identified transport channels, neutralizing any vendor-specific nomenclature discrepancies before the data enters the centralized TeraFlowSDN context database. Furthermore, the driver-level extraction mechanisms were heavily revised to explicitly capture the exact parameters required by the proposed \texttt{ParallelOpticalController}. Critical operational data—including the specific channel type, total spectral capacity, real-time slot availability, and precise physical channel-to-endpoint (port) mappings—are now mapped directly into the device's operational state. By systematically resolving and normalizing these attributes at the lowest level of the control plane (the driver level), the overarching controller is guaranteed a clean, unified, and highly accurate topological state, which is crucial for the efficient execution of the bitmap-based RSA algorithm.
\subsection{Context Service \& Database Schema Extensions}
Within the TeraFlowSDN architecture, the \texttt{ContextService} functions as the centralized repository and primary communication hub, orchestrating all persistent database transactions and cross-service data propagation. Because the device drivers were enhanced to capture highly granular, vendor-specific optical parameters, the underlying database schemas required corresponding structural extensions to accommodate and permanently store this influx of new telemetry. Consequently, the core data models within the \texttt{ContextService} were significantly augmented with novel attributes tailored explicitly to support the proposed spatial Routing and Spectrum Allocation (RSA) algorithm. Critical operational parameters—such as explicit physical endpoint indices, discrete flexible-grid frequency details, integer-encoded spectrum utilization bitmaps, and expanded hardware capability metrics—are now persistently recorded within the context database immediately following their extraction at the driver level. These newly modeled attributes form the foundational dataset required for the controller's multi-dimensional path and spectrum computations. To ensure this augmented dataset is securely and rapidly accessible across the distributed SDN control plane, the overarching Protocol Buffer (protobuf) definitions were substantially revised. By updating and recompiling these protobuf schemas, the system guarantees seamless, high-speed gRPC communication between internal components. This critical integration allows the \texttt{ParallelOpticalController} to instantly and reliably retrieve the complex spectral and spatial state of the multi-fiber network, thereby enabling real-time, mathematically accurate lightpath provisioning.
\section{Algorithmic Core and Helper Modules}
While the \texttt{ParallelOpticalController} operates as a high-level orchestration micro-service within the TeraFlowSDN ecosystem, its underlying computational intelligence—often conceptualized as the "brain" of the controller—is systematically distributed across a suite of specialized Python submodules. Rather than functioning as a monolithic script, the controller relies on these distinct internal layers working in tandem to execute complex spatial and spectral logic.

To maintain modularity and reusability, these helper functions are categorized into two domains: universally shared libraries and controller-specific core modules.
\subsection*{\small Shared Utility Modules}
To ensure strict compliance with physical layer protocols across the entire control plane, the system introduces \texttt{ITUStandards.py}. This module encapsulates the comprehensive ITU-T specifications for DWDM flexible-grid communications. During the device discovery and driver extraction phase (as detailed in Section 4.1.1), this module provides the definitive reference parameters necessary to accurately map physical hardware endpoints into standardized database formats. Complementing the standards definitions is \texttt{RSATools.py}, which provides a robust library of utility functions designed to parse and align this standardized physical data against the specific parameters of incoming lightpath requests (e.g., converting requested bitrate into required frequency slots). Because these foundational mapping and conversion operations are critical across multiple stages of the network lifecycle, both \texttt{ITUStandards.py} and \texttt{RSATools.py} are deployed as shared dependencies utilized by the \texttt{DeviceService}, the \texttt{WebService} (for accurate WebUI representation), and the \texttt{ParallelOpticalController}.
\subsection*{\small Controller-Specific Modules}
Conversely, the actual spatial routing and spectral logic is heavily abstracted into two dedicated modules strictly local to the \texttt{ParallelOpticalController}: \texttt{topology.py} and \texttt{RSAHelper.py}. These files effectively function as the internal API and computational engine of the micro-service.
\begin{itemize}
    \item \texttt{topology.py}: This module is strictly responsible for graph formulation. It handles the complex interactions with the context database to extract real-time node and link states, subsequently constructing the multi-graph and simplified logical graph models required for path discovery.
    \item \texttt{RSAHelper.py}: Acting as the primary algorithmic executor, this module handles the heavy data processing. It manages the extraction of path-specific capabilities, performs the bitmap conversions and spatial bitwise intersections across parallel links, and ultimately generates the validated spectrum allocation response.
\end{itemize}
Together, these local and shared submodules synthesize raw database telemetry into actionable intelligence, allowing the overarching \texttt{ParallelOpticalController} micro-service to efficiently orchestrate multi-fiber lightpath provisioning.
\section{Parallel Optical Controller Microservice}
Architecturally, the \texttt{ParallelOpticalController}\cite{chafiTFSGit} is deployed\footnote{The controller is not officially available and suggested to acquire from the git repository mentioned in the later section.} as an independent microservice within the TeraFlowSDN ecosystem. Unlike the mandatory core infrastructural components of the SDN, it is deliberately designed as a decoupled, portable solution. This modularity allows network operators to integrate advanced multi-fiber routing capabilities seamlessly without permanently altering the foundational control plane.
\begin{figure}[htbp]
    \centering
    \includegraphics[width=1\textwidth]{poc_service_orchestration.png}
    \caption{Service Architecture of ParallelOpticalController}
    \label{fig:poc_service_orchestration}
\end{figure}
Adhering to official TeraFlowSDN community guidelines, the controller is containerized and orchestrated via MicroK8s. To ensure full interoperability and successfully expose its advanced routing options to the overarching system, the deployment process requires a coordinated restart of specific adjacent microservices (such as the \texttt{DeviceService} and \texttt{ContextService}). A sample automation script used to execute this deployment and synchronization procedure is provided below:
\begin{lstlisting}[language=bash, caption={Deployment of ParallelOpticalController}]
    # On Primary server
    # if you face caching issue, please add --no-cache to the build command
    $ docker build -t localhost:32000/tfs/parallelopticalcontroller:dev -f \ 
    src/parallelopticalcontroller/Dockerfile .

    $ docker push localhost:32000/tfs/parallelopticalcontroller:dev
    
    # if you are building for the first time or have changes in the files names, or ports, or any thing else, update the yaml file and push it again
    $ kubectl apply -f manifests/parallelopticalcontrollerservice.yaml -n tfs
    # otherwise just restart the pod
    $ kubectl rollout restart deployment/parallelopticalcontrollerservice -n tfs
\end{lstlisting}

Once successfully deployed and integrated, the controller assumes an event-driven operational lifecycle. When a new lightpath request is initiated, the TeraFlowSDN core generates a corresponding service request object, which is subsequently exposed to the \texttt{ParallelOpticalController}. However, rather than executing autonomously in the background, the core RSA computation is explicitly triggered by the network operator via the TeraFlowSDN WebUI. This manual trigger initiates a cascade of API interactions and localized function calls across the controller's internal modules (as detailed in the previous helper functions section). The microservice retrieves the real-time context data, formulates the graph, and executes the bitmap-based spatial RSA and slot acquisition protocols. Finally, the controller aggregates the computed data—including the validated physical paths, spatial port mappings, and the final spectrum bitmaps—into a unified, structured response payload. This payload is then transmitted back to the \texttt{WebService}, translating the complex mathematical matrices into an intuitive visual representation for the operator's final review and approval.